\documentclass[11pt]{article}
\usepackage[margin=1in]{geometry}
\usepackage{amsmath,amssymb,amsthm,hyperref}
\newtheorem{lemma}{Lemma}
\newtheorem{proposition}{Proposition}
\newcommand{\tr}{\operatorname{Tr}}
\newcommand{\norm}[1]{\left\|#1\right\|}
\begin{document}
\title{Coordinate Wilson slabs with both types of reflection boundary}
\date{16 October 2026}\maketitle
\noindent Research derivation for independent audit. This note proves a
finite-volume transfer estimate. It does not prove the two-frame fold
identities, their placement, or same-limit diagonal reflection positivity.
All constants and limits in this note are at fixed positive lattice
coupling $\beta$; no uniform continuum estimate is claimed.

\section{Normalization and the physical Hilbert space}

Let $\Gamma$ be a finite three-dimensional periodic spatial lattice, with
positive links $E$ and sites $V_\Gamma$. Let $G=SU(2)$ with normalized
Haar measure. On $\mathcal H=L^2(G^E)$, the spatial gauge group
$G^{V_\Gamma}$ acts by
$(g\cdot U)_e=g_{x(e)}U_eg_{y(e)}^{-1}$.
Let $P$ be its Haar-averaged orthogonal projection, and put
$\mathcal H_{\rm phys}=P\mathcal H$.

Write
\[
 k_\beta(g)=e^{(\beta/2)\tr g},\qquad
 C=\bigotimes_{e\in E}C_\beta,
 \quad(C_\beta f)(v)=\int k_\beta(vu^{-1})f(u)\,du,
\]
and let $V$ denote multiplication by the product of $k_\beta(U_p)$ over
all spatial plaquettes. Both $C$ and $V$ commute with $P$. This follows
for $C$ from
\[
 k_\beta((g_xvg_y^{-1})(g_xug_y^{-1})^{-1})
 =k_\beta(g_xvu^{-1}g_x^{-1})=k_\beta(vu^{-1});
\]
for $V$ it is ordinary gauge invariance.
The convolution eigenvalues are positive and decrease faster than any
power of representation degree. Hence $C$ is positive, injective and
trace class, and $C^{1/2}$ is Hilbert--Schmidt. The positive function
$V$ and its square root are bounded on this finite compact space.

All following operators are restricted to $\mathcal H_{\rm phys}$.
Define
\[
 B=V^{1/2}C^{1/2},\qquad
 A_l=B^*B=C^{1/2}VC^{1/2},\qquad
 A_s=BB^*=V^{1/2}CV^{1/2}.
\]
Both $A_l,A_s$ are positive trace-class operators and have the same
nonzero eigenvalues with multiplicity. Therefore
\[
 Z_\Gamma(N)=\tr A_l^N=\tr A_s^N,\qquad N\in\mathbb N,
 \tag{1}
\]
is the ordinary periodic Wilson partition function of temporal length
$N$, for lengths avoiding degenerate one-cell plaquette identifications.
In the applications all lengths tend to infinity.

To check the normalization in (1), integrate each slice of temporal
links with normalized Haar measure. If $a_x$ is such a temporal link,
the electric plaquette trace is
$\tr(U_e a_{y(e)}U_e'{}^{-1}a_{x(e)}^{-1})$.
Thus the electric kernel, after the temporal integration, is the
kernel of $CP$. Spatial weights give $V$ once per slice. The torus
integral is consequently $\tr_{\mathcal H}(CPV)^N$, which equals
the trace on $\mathcal H_{\rm phys}$. Cycling the square roots yields
(1). There are no gauge-volume factors because all Haar measures have
mass one.

\section{Slabs and alternating boundary types}

A link-plane cut is halfway between adjacent integer time slices.
Its two boundary pieces carry $C^{1/2}$. A site-plane cut is at an
integer slice and divides that slice's spatial weight into $V^{1/2}$
on each side. The two boundary spaces can both be represented by
$\mathcal H_{\rm phys}$, with the displayed factors keeping track of
the boundary type.
From a link boundary to the next site boundary the half-step is $B$;
the reverse type of half-step is $B^*$. Consequently a slab of length
$\ell\in\tfrac12\mathbb N$ is a word of $2\ell$ alternating half-steps.
For like boundary types this gives $A_l^\ell$ or $A_s^\ell$, with
$\ell$ integral. For different types it gives a word with odd length.
The gauge projectors at temporal transitions may be retained explicitly;
on the physical space they are identities.

Let $R$ be a real bounded gauge-invariant function of links wholly
inside the slab, including the appropriate boundary-side data, but not
an interaction crossing either cut. Insert $e^{tR}$ in its Wilson path
integral and denote the resulting boundary operator by $S_{R,t}$.
If a site boundary is used, any source term reading boundary variables
is assigned to one side before the cut; terms requiring both sides
are excluded by the no-crossing assumption. In the intended application
the source boxes avoid the cut, so this assignment issue does not arise.
Assume $\norm R_\infty\leq M$.

\begin{lemma}[Even Schatten bound for a bounded slab source]
For every even integer $p\geq2$,
\[
 \norm{S_{R,t}}_p
 \leq e^{|t|M} Z_\Gamma(p\ell)^{1/p}.
 \tag{2}
\]
The integer $p\ell$ is the temporal length of the reflected closed slab
array. The two possible boundary types obey the same formula.
\end{lemma}
\begin{proof}
Reflection of the slab reverses its time ordering and takes the
operator adjoint, because all Wilson and source weights are real.
Glue $p$ alternating copies of the slab and its reflection. Adjacent
link-boundary factors multiply to $C^{1/2}C^{1/2}=C$; adjacent
site-boundary factors multiply to $V^{1/2}V^{1/2}=V$.
Temporal Haar integration implements the intervening gauge projectors.
Thus
\[
 \tr[(S_{R,t}^*S_{R,t})^{p/2}]
\]
is the ordinary positive Wilson integral on temporal length $p\ell$
with $p$ reflected source copies, each of supremum norm at most $M$.
Its value is at most $e^{p|t|M}Z_\Gamma(p\ell)$.
This proves (2), including Hilbert--Schmidt membership by taking $p=2$.
For full operator rigor one may first expand smooth boundary kernels
in finite representation sums and pass to the Hilbert--Schmidt limit;
the glued expression has the stated bounded positive integral.
No pointwise positivity of the kernel of $C^{1/2}$ is used or needed.
\end{proof}

\section{Bounded insertions between clean slabs}

Take an even number $J$ of cuts on a coordinate torus, with successive
positive slab lengths $\ell_1,\ldots,\ell_J$ and total integer length
$N=\sum_j\ell_j$. Allow either boundary type. Suppose no source term
crosses any cut, and let $R_j$ be the sum assigned to slab $j$,
$\norm{R_j}_\infty\leq M_j$.
Let $X_j$ be bounded operators on the respective physical boundary
spaces. Their supports and gauge-covariance requirements matter for
the observable interpretation, but not for the following trace bound.

\begin{proposition}[Coordinate insertion bound]
The normalized cyclic insertion trace obeys
\[
 \frac{\left|\tr(X_J S_{R_J,t}\cdots X_1S_{R_1,t})\right|}
 {Z_{\Gamma,t}(N)}
 \leq e^{2|t|\sum_jM_j}\prod_j\norm{X_j}
 \frac{\prod_j Z_\Gamma(J\ell_j)^{1/J}}{Z_\Gamma(N)}.
 \tag{3}
\]
Here the denominator is the same Wilson torus with all source terms
and no operator insertions.
\end{proposition}
\begin{proof}
Schatten H\"older with all exponents equal to $J$ bounds the numerator
by $\prod_j\norm{X_j}\norm{S_{R_j,t}}_J$.
Apply (2). Pointwise positivity of the original Wilson integral gives
$Z_{\Gamma,t}(N)\geq e^{-|t|\sum_jM_j}Z_\Gamma(N)$.
Combining the two inequalities proves (3).
\end{proof}

For a local operator $X$ on selected spatial links $S$ at a link cut,
the physical insertion used here has kernel
$C_S^{1/2}XC_S^{1/2}$ in place of $C_S$. Dividing by the strictly
positive product electric kernel gives the bounded local observable
when this numerator kernel is continuous. Finite representation
operators and the identity have this property.
If $X$ commutes with endpoint gauge actions on $S$, then $X\otimes I$
commutes with the full gauge projection. Disjoint insertions at the
same cut can be combined with norm at most the product of their norms.

\section{What the pressure limit does and does not supply}

Suppose the spatial sides, $N$ and every $J\ell_j$ tend to infinity
with bounded aspect ratios. If all these rectangular Wilson tori have
the common thermodynamic pressure $p_\beta$, then
\[
 \frac1J\sum_j\log Z_\Gamma(J\ell_j)-\log Z_\Gamma(N)
 =o(|\Gamma|N),
 \tag{4}
\]
provided the pressure errors are uniform for this family. A finite set
of fixed slab lengths, with $J$ proportional to $N$, gives precisely
such a bounded-aspect family; uniform pressure convergence follows
from the usual bounded-interaction boundary estimate.
The leading terms cancel because
$J^{-1}\sum_j p_\beta|\Gamma|J\ell_j=p_\beta|\Gamma|N$.

If a separately proved chessboard amplification creates $c|\Gamma|N$
copies of one local operator, its root converts the exponential of
(4) to $e^{o(1)}$ and the bounded-source term in (3) to $e^{O(t)}$.
This is the desired analytic end of the amplification calculation.
The required two-frame fold identity, number of copies, no-crossing
placement and transport of local insertions must still be established
for the actual adapted-torus geometry. Formula (3) by itself does not
transfer reflection positivity between two limiting laws.

\section{The exact consequence of a completed amplification estimate}

Suppose for one coordinate link plane and every local finite
representation-kernel operator $X$, together with the identity, a
comparison limit $\lambda$ obeys
\[
 |\lambda(G_X)|\leq\norm X,\qquad G_I=1.
 \tag{5}
\]
Here $G_X$ is the gauge-invariant crossing insertion defined above,
and the inequality is also valid for $I+itA$ with self-adjoint $A$.
Set $L(X)=\lambda(G_X)$. For self-adjoint $A$, (5) gives
\[
 |1+itL(A)|^2\leq\norm{I+itA}^2
                  \leq1+t^2\norm A^2.
\]
Divide the resulting inequality by $t$ and let $t\to0$ through both
signs. It follows that $\operatorname{Im}L(A)=0$, and hence
$L(X^*)=\overline{L(X)}$ for all such $X$.
Reflection swaps the crossing kernel arguments, so
$G_X\circ\theta=\overline{G_{X^*}}$ after the temporal parallel
transports are included. Therefore
\[
 \lambda(G_X\circ\theta)=\lambda(G_X).
 \tag{6}
\]

Every smooth gauge-invariant two-layer observable $Y$ is uniformly
approximated by these insertions. In transported temporal coordinates,
let $K_Y(v,u)=K_S(v,u)Y(u,v)$, and let $P_m$ be the product
Peter--Weyl cutoff on the selected finite edge set. The smooth kernels
$P_mK_YP_m$ converge uniformly to $K_Y$. Define
\[
 X_m=C_S^{-1/2}P_mK_YP_mC_S^{-1/2}.
\]
This is a bounded finite-rank operator, since every retained
convolution eigenvalue is strictly positive. It commutes with the
endpoint gauge action, since all the factors do. Its insertion is
$G_{X_m}=P_mK_YP_m/K_S$, so it converges uniformly to $Y$ because
$K_S$ has a positive lower bound on the finite compact link space.
Consequently (6) holds for all continuous gauge-invariant two-layer
observables by a further uniform approximation.
No bound uniform in $m$ on $\norm{X_m}$ is required: first remove
the tilt and obtain (5) for each fixed operator, then approximate $Y$.

Finally, suppose the paired two-layer reduction has been proved for
each local $P$ and its reflection:
\[
 |\lambda(P)-\lambda(Y_\epsilon)|\leq\epsilon,\qquad
 |\lambda(P\circ\theta)-\lambda(Y_\epsilon\circ\theta)|\leq\epsilon.
\]
Equation (6) and $\epsilon\downarrow0$ give reflection invariance for
all local observables. This proves invariance, not reflection
positivity at that new coordinate plane. In the intended convex
separation argument it is used to obtain even translations by composing
two such invariances; diagonal positivity comes from the original
adapted-torus construction.
\end{document}
